Many materials and tissues, including carbon-nanotube sheets, collagen, paper
and nonwoven textiles, comprise overlapping strands whose arrangement
influences their effective properties. Quantifying this geometry from
micrographs requires individual strands to be recovered through crossings and
signal gaps, which merge distinct strands or split one strand into disconnected
fragments. We present PLECTA, a deterministic method that reconstructs
individual strands using only a binary image whose foreground pixels mark
strand centrelines. At each junction, PLECTA selects which strand segments
continue into one another while allowing segments to terminate. It formulates
this decision as an exact maximum-weight partial-matching problem, with an
explicit penalty for unpaired segments and no restriction on the number of
segments entering a junction. The output permits crossing pixels to belong to
multiple strands. On procedurally generated scenes, PLECTA reconstructed
strand identities more accurately than three released alternatives across all
evaluated densities, and this ranking was preserved on an independently
published generator evaluated using its original settings. A separate stage
uses the corresponding greyscale micrograph to estimate which strand lies on
top at each crossing, providing a proof of concept for depth ordering. The
complete workflow is demonstrated on scanning electron micrographs of
carbon-nanotube sheets, from raw images to individual strand instances with
estimated crossing order, without manual annotations for training or as
pipeline inputs.
